% Wersja źródłowa: audyt II, jednostka 29
\setcounter{section}{28}
\section{Operatory ograniczone i zwarte: kryteria, transpozycja i ideały}
\label{sec:operatory29}
\noindent\textbf{Wyniki wykorzystywane.}
Korzystamy z normującej postaci twierdzenia Hahna--Banacha
(Wniosek~\ref{cor:norming}), zasady jednostajnej ograniczoności
(Wniosek~\ref{cor:uniform-boundedness}), twierdzenia o odwzorowaniu
otwartym (Twierdzenie~\ref{thm:open-mapping11}), twierdzenia
Eberleina--\v Smuliana (Twierdzenie~\ref{thm:eberlein-smulian26})
oraz Twierdzenia~\ref{thm:approximable-compact28}.
Korzystamy również z klasycznych twierdzeń Arzeli--Ascolego
i Fubiniego oraz regularności miary Lebesgue'a.
\par\medskip

W rozdziale przyjmujemy, że $X,Y,Z$ są przestrzeniami Banacha nad
$\K\in\{\R,\C\}$. Zapis $\mathcal L(X,Y)$ oznacza przestrzeń
ograniczonych operatorów liniowych, $\mathcal F(X,Y)$ -- operatory
ciągłe skończonego rzędu, a $\mathcal K(X,Y)$ -- operatory zwarte.
Te pojęcia, a także domkniętość $\mathcal K(X,Y)$ w normie operatorowej,
ustanowiono w rozdziale o własnościach aproksymacji. Nie powtarzamy
tamtych dowodów. Dla operatora $T:X\to Y$ przez $T':Y'\to X'$
oznaczamy jego operator sprzężony (transpozycję), dany wzorem
$(T'y')(x)=y'(Tx)$. Normy operatora i jego transpozycji są równe:
$\|T'\|=\|T\|$. Istotnie, z Wniosku~\ref{cor:norming}
i definicji normy operatorowej wynika
\[
\|T'\|=\sup_{\|y'\|\le1}\sup_{\|x\|\le1}|y'(Tx)|
=\sup_{\|x\|\le1}\|Tx\|=\|T\|.
\]

\subsection{Zwartość i kryteria ciągowe}

\begin{twierdzenie}[Trzy kryteria zwartości]
\label{thm:compact-criteria29}
Dla $T\in\mathcal L(X,Y)$ następujące warunki są równoważne:
\begin{enumerate}[label=\textup{(\roman*)}]
\item $T$ jest zwarty, tj. $\overline{T(B_X)}$ jest zwarty w $Y$;
\item dla każdego ograniczonego ciągu $(x_n)$ w $X$ ciąg $(Tx_n)$
      ma podciąg zbieżny w normie $Y$;
\item dla każdego $\varepsilon>0$ zbiór $T(B_X)$ ma skończoną
      $\varepsilon$-siatkę w $Y$.
\end{enumerate}
\end{twierdzenie}
\begin{proof}
(i) $\Rightarrow$ (ii): dowolny ograniczony ciąg $(x_n)$
mieści się w $rB_X$ dla pewnego $r>0$. Zbiór
$r\overline{T(B_X)}$ jest zwarty, więc ciąg $(Tx_n)$ ma
podciąg zbieżny w normie.
(ii) $\Rightarrow$ (iii): gdyby $T(B_X)$ nie miał skończonej
$\varepsilon$-siatki dla pewnego $\varepsilon>0$, wybieralibyśmy
indukcyjnie $x_n\in B_X$ tak, aby $Tx_{n+1}$ leżał poza sumą
kul otwartych promienia $\varepsilon$ wokół $Tx_1,\ldots,Tx_n$.
Wtedy $\|Tx_n-Tx_m\|\ge\varepsilon$ dla $n\ne m$,
więc żaden podciąg nie byłby ciągiem Cauchy'ego.
(iii) $\Rightarrow$ (i): warunek (iii) oznacza całkowitą
ograniczoność $T(B_X)$. Domknięcie tego zbioru jest również
całkowicie ograniczone: dla zadanego promienia wybieramy
siatkę o połowie tego promienia, a każdy punkt domknięcia
przybliżamy punktem zbioru z błędem mniejszym niż połowa
promienia. Domknięty podzbiór zupełnej przestrzeni $Y$ jest
zupełny. Z metrycznego kryterium zwartości otrzymujemy (i).
\end{proof}

\begin{lemat}[Lemat Riesza o prawie prostopadłym wektorze]
\label{lem:riesz29}
Niech $M$ będzie właściwą domkniętą podprzestrzenią przestrzeni
unormowanej $X$, a $0<\theta<1$. Istnieje $x\in X$ takie, że
$\|x\|=1$ oraz $\operatorname{dist}(x,M)>\theta$.
\end{lemat}
\begin{proof}
Wybieramy $z\in X\setminus M$. Ponieważ $M$ jest domknięta,
$d:=\operatorname{dist}(z,M)>0$. Wybieramy $y\in M$ z
$\|z-y\|<d/\theta$ i kładziemy $x=(z-y)/\|z-y\|$.
Dla każdego $m\in M$ zachodzi
$\|x-m\|=\|z-y-\|z-y\|m\|/\|z-y\|\ge d/\|z-y\|$,
ponieważ $y+\|z-y\|m\in M$. Biorąc infimum po $m\in M$,
dostajemy $\operatorname{dist}(x,M)\ge d/\|z-y\|>\theta$.
\end{proof}

\begin{twierdzenie}[Zwartość identyczności]
\label{thm:identity-compact29}
Dla przestrzeni Banacha $X$ operator $I_X$ jest zwarty wtedy
i tylko wtedy, gdy $X$ ma skończony wymiar.
\end{twierdzenie}
\begin{proof}
W skończonym wymiarze ograniczone zbiory mają zwarte domknięcie.
Gdy $X$ ma nieskończony wymiar, wybieramy indukcyjnie wektory
$x_n\in B_X$ o normie $1$ w taki sposób, że
$\operatorname{dist}(x_{n+1},\operatorname{span}\{x_1,\ldots,x_n\})>1/2$.
Lemat~\ref{lem:riesz29} pozwala wykonać każdy krok, gdyż
podprzestrzeń skończenie wymiarowa jest domknięta i właściwa.
Wówczas $\|x_n-x_m\|>1/2$ dla $n\ne m$, więc kula jednostkowa
nie jest względnie zwarta.
\end{proof}

\begin{twierdzenie}[Zwartość a składanie operatorów]
\label{thm:ideal-compact29}
Jeżeli $S\in\mathcal L(Y,Z)$ i $T\in\mathcal L(X,Y)$, to:
\begin{enumerate}[label=\textup{(\roman*)}]
\item $T\in\mathcal K(X,Y)$ implikuje
      $ST\in\mathcal K(X,Z)$;
\item $S\in\mathcal K(Y,Z)$ implikuje
      $ST\in\mathcal K(X,Z)$.
\end{enumerate}
Ponadto $\mathcal K(X,Y)$ jest podprzestrzenią liniową zamkniętą
w normie operatorowej.
\end{twierdzenie}
\begin{proof}
W (i) zbiór $T(B_X)$ jest względnie zwarty, a ciągłość $S$
zapewnia względną zwartość jego obrazu. W (ii) zbiór $T(B_X)$
jest ograniczony, więc zawiera się w $rB_Y$ dla pewnego $r>0$;
$S(rB_Y)$ jest względnie zwarty. Ostatnie zdanie jest treścią
Twierdzenia~\ref{thm:approximable-compact28}.
\end{proof}

\begin{przyklad}[Zwartość nie wymaga skończonego rzędu]
\label{ex:volterra29}
W przestrzeni Banacha $C([0,1],\K)$ definiujemy operator Volterry
$(Vf)(t)=\int_{0}^{t}f(s)\,\mathrm d\lambda_1(s)$.
Liniowość wynika z liniowości całki, a oszacowanie
$|(Vf)(t)|\le t\|f\|_\infty$ daje $\|V\|\le1$.
Wartości funkcji $Vf$ zależą od $t$ w sposób ciągły, gdyż
$|(Vf)(t)-(Vf)(u)|\le\|f\|_\infty|t-u|$.
Dla $\|f\|_\infty\le1$ mamy $\|Vf\|_\infty\le1$ i
$|(Vf)(t)-(Vf)(u)|\le|t-u|$. Zatem rodzina $V(B_X)$
jest jednostajnie ograniczona i jednakowo ciągła;
z twierdzenia Arzeli--Ascolego ma zwarte domknięcie.
Operator $V$ jest zatem zwarty. Jednocześnie obrazy wielomianów
$f_n(t)=t^n$, równe $t^{n+1}/(n+1)$, są liniowo niezależne,
więc rząd $V$ jest nieskończony. Ponadto z podstawowego twierdzenia rachunku całkowego
$(Vf)'=f$ na $(0,1)$, więc $Vf=0$ implikuje $f=0$
także na końcach odcinka przez ciągłość. Operator jest
różnowartościowy, ale jego obraz nie jest domknięty: gdyby był domknięty, odwrotność
$V^{-1}:V(X)\to X$ byłaby ciągła z twierdzenia o odwzorowaniu
otwartym, a stąd $I_X=V^{-1}V$ byłaby zwarta, wbrew
Twierdzeniu~\ref{thm:identity-compact29}.
\end{przyklad}

\subsection{Twierdzenie Schaudera o operatorze sprzężonym}

\begin{lemat}[Jednostajna aproksymacja ograniczonych funkcjonałów]
\label{lem:restrictions29}
Niech $K$ będzie zwartym podzbiorem przestrzeni Banacha $Y$.
Rodzina $\{y'|_K:\|y'\|\le1\}$ jest względnie zwarta
w $C(K,\K)$ z normą supremum.
\end{lemat}
\begin{proof}
Dla $k\in K$ mamy $|y'(k)|\le\|k\|$, więc rodzina jest
jednostajnie ograniczona. Ponadto
$|y'(k)-y'(l)|\le\|k-l\|$ dla dowolnych $k,l\in K$:
rodzina jest jednakowo ciągła. Twierdzenie Arzeli--Ascolego
(dla funkcji zespolonych stosujemy je do części rzeczywistych
i urojonych) daje względną zwartość w normie supremum.
\end{proof}

\begin{twierdzenie}[Schauder: zwartość transpozycji]
\label{thm:schauder-adjoint29}
Niech $X,Y$ będą przestrzeniami Banacha nad tym samym ciałem
$\K$, a $T\in\mathcal L(X,Y)$. Wówczas
\[
T\in\mathcal K(X,Y)\quad\Longleftrightarrow\quad
T'\in\mathcal K(Y',X').
\]
\end{twierdzenie}
\begin{proof}
Załóżmy najpierw, że $T$ jest zwarty, i niech
$K=\overline{T(B_X)}\subseteq Y$. Dla $y'_1,y'_2\in B_{Y'}$
z definicji transpozycji dostajemy
\[
\|T'y'_1-T'y'_2\|_{X'}
=\sup_{x\in B_X}|(y'_1-y'_2)(Tx)|
\le\sup_{k\in K}|(y'_1-y'_2)(k)|.
\]
Lemat~\ref{lem:restrictions29} pokazuje, że rodzina
restrykcji $y'|_K$ ma zwarte domknięcie w $C(K,\K)$.
Dla każdego $\varepsilon>0$ wybieramy skończoną
$\varepsilon/3$-siatkę tej rodziny. Z każdej kuli siatki,
która przecina rodzinę, wybieramy jedną restrykcję
$y'_j|_K$ dla $y'_j\in B_{Y'}$. Każda restrykcja jest
od wybranej restrykcji odległa o mniej niż $2\varepsilon/3$.
Z powyższego oszacowania wynika, że
$T'y'_1,\ldots,T'y'_m$ tworzą skończoną $\varepsilon$-siatkę
$T'(B_{Y'})$. Twierdzenie~\ref{thm:compact-criteria29}
daje zwartość $T'$.

Jeżeli natomiast $T'$ jest zwarty, to właśnie udowodniona
implikacja, zastosowana do $T':Y'\to X'$, daje zwartość
$T'':X''\to Y''$. Oznaczmy przez $J_X:X\to X''$ i
$J_Y:Y\to Y''$ kanoniczne izometryczne zanurzenia.
Dla każdego $x\in X$ i $y'\in Y'$ mamy
$(T''J_Xx)(y')=(J_Xx)(T'y')=(T'y')(x)=y'(Tx)=(J_YTx)(y')$,
co dowodzi tożsamości $T''J_Xx=J_YTx$. Ponieważ $J_X(B_X)\subseteq B_{X''}$,
zbiór $J_YT(B_X)$ jest względnie zwarty w $Y''$.
Obraz $J_Y(Y)$ jest domknięty (ponieważ $Y$ jest zupełna),
a $J_Y:Y\to J_Y(Y)$ jest izometrią. Stąd $T(B_X)$
jest względnie zwarty w $Y$, co dowodzi tezy.
\end{proof}

\begin{wniosek}[Sprzężenie zachowuje zwartość]
\label{cor:schauder-iterated29}
Dla każdego $T\in\mathcal L(X,Y)$ zwartość $T$ jest
równoważna zwartości $T'$ i $T''$; nie wymagamy
refleksywności żadnej z przestrzeni.
\end{wniosek}
\begin{proof}
Pierwsza równoważność jest treścią
Twierdzenia~\ref{thm:schauder-adjoint29}; druga wynika
z zastosowania tego twierdzenia do $T'$ oraz z tożsamości
$T''J_X=J_YT$.
\end{proof}

\subsection{Zbieżność słaba, operatory całkowicie ciągłe i słabo zwarte}

\begin{twierdzenie}[Operator zwarty na ciągu słabo zbieżnym]
\label{thm:weak-norm29}
Niech $T\in\mathcal K(X,Y)$. Jeżeli ciąg $(x_n)$ zbiega
słabo do $x\in X$, to $\|Tx_n-Tx\|\to0$.
To samo zachodzi dla \emph{ograniczonej} sieci słabo zbieżnej.
\end{twierdzenie}
\begin{proof}
Dla ciągu słabo zbieżnego twierdzenie Banacha--Steinhausa
zastosowane do obrazów $J_Xx_n\in X''$ zapewnia
$\sup_n\|x_n\|<\infty$. Jeżeli $\|Tx_n-Tx\|$ nie zbiega
do zera, można wybrać podciąg, którego odległość od $Tx$
jest ograniczona od dołu. Ze zwartości $T$ mamy podpodciąg
zbieżny w normie do pewnego $y\in Y$. Ponieważ $T$
jest ciągły względem topologii słabych (dla $y'\in Y'$
mamy $y'(Tx_n)=(T'y')(x_n)\to (T'y')(x)$),
granica słaba tego podpodciągu
musi być równa $Tx$, więc $y=Tx$ -- sprzeczność.
Dla ograniczonej sieci jej obraz jest względnie zwarty
w normie. Gdyby obrazy nie zbiegały do $Tx$, istniałaby
kula $B(Tx,\varepsilon)$, poza którą obrazy leżą
na współkońcowym zbiorze indeksów. Możemy wówczas
wybrać podsieć o wszystkich wyrazach poza tą kulą.
Ze zwartości można wybrać dalszą podsieć normowo zbieżną,
a jej granica słaba musiałaby być równa $Tx$.
Jednak dopełnienie kuli $B(Tx,\varepsilon)$ jest normowo
domknięte, więc granica normowa należy do tego dopełnienia.
Otrzymujemy sprzeczność.
\end{proof}

\begin{definicja}[Operator całkowicie ciągły i słabo zwarty]
\label{def:completely-continuous29}
Operator $T\in\mathcal L(X,Y)$ nazywamy
\emph{całkowicie ciągłym}, jeżeli odwzorowuje każdy ciąg
słabo zbieżny w $X$ na ciąg zbieżny w normie $Y$.
Operator $T$ jest \emph{słabo zwarty}, jeżeli
$T(B_X)$ ma słabo zwarte domknięcie w $Y$.
\end{definicja}

\begin{przyklad}[Implikacje są ostre]
\label{ex:weak-vs-compact29}
Każdy operator zwarty jest całkowicie ciągły na mocy
Twierdzenia~\ref{thm:weak-norm29} i słabo zwarty, ponieważ
normowo zwarte domknięcie obrazu kuli jest słabo zwarte
oraz słabo domknięte (topologia słaba jest Hausdorffa).
Jednak operator identyczności na $\ell^1$ jest całkowicie
ciągły na mocy własności Schura, choć nie jest zwarty,
ponieważ $\ell^1$ ma nieskończony wymiar. Nie jest też
słabo zwarty: ciąg $(e_n)$ nie ma podciągu słabo zbieżnego,
gdyż własność Schura (Twierdzenie~\ref{thm:schur-l1-26})
wymuszałaby zbieżność normową, a $\|e_n-e_m\|_1=2$
dla $n\ne m$. Stosujemy twierdzenie Eberleina--\v Smuliana.
Z drugiej strony naturalne zanurzenie
$J:\ell^1\to\ell^2$, $J(x)=x$, jest ograniczone
($\|x\|_2\le\|x\|_1$) i słabo zwarte, gdyż obraz kuli
leży w słabo zwartej kuli refleksywnej przestrzeni $\ell^2$.
Jest także całkowicie ciągłe dzięki własności Schura,
ale nie jest zwarte: $Je_n=e_n$ i $\|e_n-e_m\|_2=\sqrt2$
dla $n\ne m$. Natomiast $I_{\ell^2}$ jest słabo zwarty,
lecz nie jest całkowicie ciągły, ponieważ $e_n\rightharpoonup0$
(współczynniki każdego wektora z $\ell^2$ dążą do zera),
ale $\|e_n\|_2=1$.
\end{przyklad}

\begin{twierdzenie}[Kryterium w refleksywnej przestrzeni Banacha]
\label{thm:reflexive-complete-cont29}
Jeżeli $X$ jest refleksywną przestrzenią Banacha, to
$T\in\mathcal L(X,Y)$ jest zwarty wtedy i tylko wtedy,
gdy jest całkowicie ciągły.
\end{twierdzenie}
\begin{proof}
Jedna implikacja wynika z
Twierdzenia~\ref{thm:weak-norm29}. Odwrotnie,
niech $(x_n)$ będzie dowolnym ciągiem z $B_X$.
Refleksywność daje słabą zwartość $B_X$, a twierdzenie
Eberleina--\v Smuliana zapewnia podciąg $(x_{n_k})$
słabo zbieżny do pewnego $x\in B_X$.
Całkowita ciągłość $T$ daje normową zbieżność
$Tx_{n_k}\to Tx$. Stosujemy
Twierdzenie~\ref{thm:compact-criteria29}.
\end{proof}

\subsection{Aproksymacja operatorów zwartych i przykłady funkcyjne}

\begin{twierdzenie}[Kryterium ogonów w przestrzeniach ciągowych]
\label{thm:tails29}
Niech $Y=c_0$ lub $Y=\ell^p$ dla $1\le p<\infty$,
a $P_N:Y\to Y$ oznacza obcięcie ciągu po $N$ współrzędnych.
Dla każdego $T\in\mathcal L(X,Y)$ następujące warunki są
równoważne:
\begin{enumerate}[label=\textup{(\roman*)}]
\item $T$ jest zwarty;
\item $\lim_{N\to\infty}\|(I_Y-P_N)T\|=0$;
\item $\lim_{N\to\infty}\sup_{x\in B_X}
      \|(I_Y-P_N)Tx\|=0$.
\end{enumerate}
\end{twierdzenie}
\begin{proof}
Równoważność (ii) i (iii) jest konsekwencją definicji normy
operatorowej, gdyż dla każdego $N$ norma operatora
$(I_Y-P_N)T$ jest równa supremum w (iii).
Niech $T$ będzie zwarty i $K=\overline{T(B_X)}$.
Ponieważ $P_Ny\to y$ dla każdego $y\in Y$, a
$\|I_Y-P_N\|\le1$, zbieżność jest jednostajna na $K$:
wybieramy skończoną siatkę w $K$ i stosujemy oszacowanie
$\|(I-P_N)y\|\le\|(I-P_N)(y-y_j)\|+
\|(I-P_N)y_j\|$. Dla zadanej $\varepsilon>0$ wybieramy
siatkę promienia $\varepsilon/3$; dla dostatecznie dużych $N$
wszystkie skończenie liczne drugie składniki są mniejsze
niż $\varepsilon/3$. Supremum po $K$ jest więc nie większe
niż $2\varepsilon/3$. Stąd (ii).
W drugą stronę $P_NT$ jest operatorem skończonego rzędu,
a więc jest zwarty. Z założenia $P_NT\to T$ w normie
operatorowej, a domkniętość klasy operatorów zwartych
daje (i).
\end{proof}

\begin{przyklad}[Operatory diagonalne]
\label{ex:diagonal29}
Niech $Y=c_0$ albo $Y=\ell^p$, $1\le p<\infty$,
oraz $a=(a_n)\in\ell^\infty$. Operator
$D_a:Y\to Y$, $(D_ax)_n=a_nx_n$, jest ograniczony
i $\|D_a\|\le\|a\|_\infty$ z oszacowania współrzędnych.
Ponieważ $\|D_ae_n\|=|a_n|$ i $\|e_n\|=1$, nierówność
odwrotna wynika po wzięciu supremum po $n$. Twierdzenie~\ref{thm:tails29}
pokazuje, że $D_a$ jest zwarty dokładnie wtedy, gdy
$a_n\to0$, ponieważ
$\|(I-P_N)D_a\|=\sup_{n>N}|a_n|$.
Na przykład $D_a$ dla $a_n=1/n$ jest zwartym
operatorem różnowartościowym o niedomkniętym obrazie:
gdyby obraz był domknięty, twierdzenie o odwzorowaniu
otwartym dawałoby ciągłość odwrotności na obrazie, a zatem
$\|e_n\|\le C\|D_ae_n\|=C/n$, co jest niemożliwe.
\end{przyklad}

\begin{twierdzenie}[Operatory całkowe o jądrze kwadratowalnym]
\label{thm:hilbert-schmidt29}
Niech $k\in L^2([0,1]^2,\lambda_2)$ i niech
$H=L^2([0,1],\lambda_1)$ będzie przestrzenią Hilberta.
Wówczas wzór
\[
(T_kf)(t)=\int_{[0,1]}k(t,s)f(s)\,\mathrm d\lambda_1(s)
\]
definiuje dla każdego $f\in H$ funkcję prawie wszędzie
określoną, której klasa w $H$ nie zależy od reprezentantów
$k$ i $f$. Otrzymujemy operator liniowy ciągły $T_k:H\to H$. Jest on zwarty i spełnia
$\|T_k\|\le\|k\|_{L^2(\lambda_2)}$.
\end{twierdzenie}
\begin{proof}
Z twierdzenia Fubiniego dla $|k|^2$ wynika, że
$k(t,\cdot)\in L^2$ dla prawie każdego $t$. Dla $f\in H$
iloczyn $k(t,s)f(s)$ jest całkowalny na kwadracie dzięki
nierówności Cauchy'ego--Schwarza w $L^2(\lambda_2)$,
ponieważ funkcja $(t,s)\mapsto f(s)$ ma normę $\|f\|_2$.
Fubini zapewnia mierzalność całki względem $s$ oraz
niezależność jej klasy od zmian reprezentantów na zbiorach
miary zero. Liniowość wynika z liniowości całki.
Dla prawie każdego $t$ nierówność Cauchy'ego--Schwarza daje:
$|(T_kf)(t)|^2\le
(\int_{[0,1]}|k(t,s)|^2\,\mathrm d\lambda_1(s))\|f\|_2^2$.
Całkując po $t$ względem $\lambda_1$ otrzymujemy
$\|T_kf\|_2\le\|k\|_2\|f\|_2$.
Aby uzasadnić potrzebną aproksymację jądra, zauważmy, że
funkcje charakterystyczne prostokątów $A\times B$ mają postać
$\chi_A(t)\chi_B(s)$, a ich skończone kombinacje
liniowe są gęste w $L^2([0,1]^2,\lambda_2)$.
Wynika to z gęstości funkcji prostych oraz aproksymacji
zbiorów mierzalnych przez skończone sumy prostokątów
w pseudometryce $A,B\mapsto\lambda_2(A\triangle B)$
(identyfikującej zbiory różniące się o zbiór miary zero).
Istotnie, z regularności miary Lebesgue'a każdy mierzalny
zbiór $E$ w kwadracie można zawrzeć pomiędzy zbiorem
zwartym $C$ i względnie otwartym $O$ z
$\lambda_2(O\setminus C)<\delta$. Każdy punkt $C$ leży
w względnie otwartym prostokącie zawartym w $O$.
Ze zwartości wybieramy skończone pokrycie $C$ takimi
prostokątami; ich suma $R$ spełnia
$\lambda_2(E\triangle R)<\delta$.
Funkcję $\chi_R$ zapisujemy jako skończoną kombinację
funkcji charakterystycznych prostokątów przez zasadę
włączeń i wyłączeń. Ponieważ
$\|\chi_E-\chi_R\|_2^2=\lambda_2(E\triangle R)$,
gęstość funkcji prostych kończy uzasadnienie.
Wybieramy zatem jądra $k_n$ będące skończonymi sumami
$g_j(t)h_j(s)$ i zbiegające do $k$ w $L^2$.
Jeżeli $k_n(t,s)=\sum_{j=1}^{m_n}g_j(t)h_j(s)$,
gdzie $g_j,h_j\in L^2([0,1],\lambda_1)$, to
\[
(T_{k_n}f)(t)=\sum_{j=1}^{m_n}g_j(t)
\int_{[0,1]}h_j(s)f(s)\,\mathrm d\lambda_1(s).
\]
Każdy współczynnik jest dobrze określony na mocy nierówności
Cauchy'ego--Schwarza. Obraz $T_{k_n}$ zawiera się więc
w $\operatorname{span}\{g_1,\ldots,g_{m_n}\}$. Z oszacowania normy wynika
$\|T_k-T_{k_n}\|\le\|k-k_n\|_2\to0$.
Zatem $T_k$ jest zwarty jako granica normowa operatorów
skończonego rzędu.
\end{proof}

\begin{uwaga}[Nie wszystkie operatory zwarte są aproksymowalne]
\label{rem:nonapproximable29}
W Twierdzeniu~\ref{thm:ap-compact-operators28} wykazano,
że jeżeli przestrzeń wartości $Y$ ma AP, to $\mathcal K(X,Y)=\mathcal A(X,Y)$.
Bez założenia AP w przestrzeni wartości takie twierdzenie
nie zachodzi automatycznie. Dlatego aproksymacja przez
operator skończonego rzędu nie jest równoważną
\emph{definicją} zwartości w dowolnych przestrzeniach Banacha.
\end{uwaga}

\subsection{Ideały operatorowe i przygotowanie teorii Fredholma}

\begin{definicja}[Ideał operatorowy]
\label{def:operator-ideal29}
Przez \emph{ideał operatorowy} w klasie przestrzeni Banacha
rozumiemy przyporządkowanie każdej parze $(X,Y)$
podprzestrzeni liniowej $\mathcal I(X,Y)\subseteq\mathcal L(X,Y)$,
zawierającej wszystkie operatory skończonego rzędu,
takie, że dla wszystkich przestrzeni Banacha $W,Z$
i operatorów $R\in\mathcal L(W,X)$, $T\in\mathcal I(X,Y)$,
$S\in\mathcal L(Y,Z)$ zachodzi
$STR\in\mathcal I(W,Z)$.
\end{definicja}

\begin{twierdzenie}[Trzy ideały i domkniętość]
\label{thm:ideals29}
Rodziny $\mathcal F$, $\mathcal A$ i $\mathcal K$ są
ideałami operatorowymi. Przy tym $\mathcal A(X,Y)$
i $\mathcal K(X,Y)$ są domknięte w normie operatorowej,
a $\mathcal F(X,Y)$ nie jest domknięta, jeżeli
$X=Y=\ell^2$. Dla wszystkich $X,Y$ zachodzi
$\mathcal F(X,Y)\subseteq\mathcal A(X,Y)\subseteq
\mathcal K(X,Y)$.
\end{twierdzenie}
\begin{proof}
Suma dwóch operatorów skończonego rzędu ma obraz zawarty
w sumie ich skończenie wymiarowych obrazów, a mnożenie
przez skalar zachowuje skończony rząd. Zatem $\mathcal F(X,Y)$
jest podprzestrzenią liniową, podobnie jej normowe domknięcie
$\mathcal A(X,Y)$. Składanie operatora skończonego rzędu z operatorami
ograniczonymi z lewej i prawej strony zachowuje
skończoność wymiaru obrazu; dlatego $\mathcal F$ jest ideałem.
Jeżeli $T_n\to T$ w normie i każdy $T_n$ ma skończony
rząd, to $ST_nR\to STR$ w normie na mocy
$\|ST_nR-STR\|\le\|S\|\|T_n-T\|\|R\|$.
To daje własność ideału $\mathcal A$ i jego domkniętość
z definicji. Twierdzenie~\ref{thm:ideal-compact29}
daje własność ideału $\mathcal K$; jego domkniętość
udowodniono wcześniej. Operator diagonalny
$D_{(1/n)}:\ell^2\to\ell^2$ ma nieskończony rząd,
a jego obcięcia skończonego rzędu zbiegają do niego
w normie. To pokazuje, że $\mathcal F(\ell^2,\ell^2)$
nie jest domknięta. Zawierania między trzema ideałami
wynikają z definicji i zwartości operatorów
skończonego rzędu.
\end{proof}

\begin{definicja}[Norma istotna operatora]
\label{def:essential-norm29}
Dla $T\in\mathcal L(X,Y)$ definiujemy
$\|T\|_{\mathrm{ess}}=
\inf\{\|T-K\|:K\in\mathcal K(X,Y)\}$.
Jest to odległość operatora $T$ od domkniętej
podprzestrzeni operatorów zwartych; może wynosić zero
dla operatora niezerowego. Nie jest więc normą
na całej przestrzeni $\mathcal L(X,Y)$, ale wyznacza normę
na przestrzeni ilorazowej
$\mathcal L(X,Y)/\mathcal K(X,Y)$. Gdy $X=Y$,
przestrzeń ta jest dodatkowo algebrą ilorazową.
\end{definicja}

\begin{twierdzenie}[Norma istotna identyczności]
\label{thm:essential-identity29}
Jeżeli $X$ jest nieskończenie wymiarową przestrzenią
Banacha, to $\|I_X\|_{\mathrm{ess}}=1$.
\end{twierdzenie}
\begin{proof}
Wybór $K=0$ pokazuje, że odległość jest nie większa
niż $1$. Załóżmy, że istnieje $K\in\mathcal K(X,X)$
z $\|I_X-K\|<1$. Stosując Lemat~\ref{lem:neumann15}
do operatora $T=I_X-K$, otrzymujemy odwracalność
$I_X-T=K$. W szczególności
$I_X=K^{-1}K$ byłby zwarty na mocy
Twierdzenia~\ref{thm:ideal-compact29}, wbrew
Twierdzeniu~\ref{thm:identity-compact29}.
\end{proof}

\begin{lemat}[Pierwszy krok teorii Fredholma]
\label{lem:compact-eigenvalue29}
Niech $T\in\mathcal K(X,X)$ i $\lambda\in\K\setminus\{0\}$.
Jeśli operator $\lambda I_X-T$ jest różnowartościowy,
to istnieje $c>0$ takie, że
$\|(\lambda I_X-T)x\|\ge c\|x\|$ dla wszystkich $x\in X$.
W szczególności obraz $\operatorname{ran}(\lambda I_X-T)$
jest domknięty.
\end{lemat}
\begin{proof}
Gdyby nie było takiego $c$, istniałby ciąg $(x_n)$
o normach równych $1$, spełniający
$(\lambda I_X-T)x_n\to0$. Ze zwartości $T$ możemy
wybrać podciąg, dla którego $Tx_{n_k}\to y$ w normie.
Wówczas $\lambda x_{n_k}=Tx_{n_k}+
(\lambda I_X-T)x_{n_k}\to y$, więc $x_{n_k}\to x=y/\lambda$
i $\|x\|=1$. Przez ciągłość operatora otrzymujemy
$(\lambda I_X-T)x=0$, co przeczy jego różnowartościowości.
Jeżeli obrazy $(\lambda I_X-T)x_n$ tworzą ciąg Cauchy'ego,
to uzyskane oszacowanie implikuje, że $(x_n)$ jest
ciągiem Cauchy'ego. Zupełność $X$ i ciągłość operatora
pokazują, że obraz jest domknięty.
\end{proof}

\begin{uwaga}[Granica zakresu rozdziału]
\label{rem:fredholm-next29}
Ostatni lemat nie stwierdza jeszcze surjektywności
$\lambda I_X-T$ ani nie opisuje krotności wartości
własnych. W następnym rozdziale dowiedziemy twierdzenia
Riesza--Schaudera, alternatywy Fredholma i stabilności
indeksu przy zaburzeniach zwartych, nie powtarzając
kryteriów zwartości i twierdzenia o transpozycji.
\end{uwaga}
